\documentclass[%
sor,
 jor,
 amsmath,amssymb,
 reprint,
longbibliography,
]{revtex4-2}
\usepackage{float}
\usepackage{caption}
\usepackage{empheq, amssymb, amsfonts}
\usepackage{siunitx}
\usepackage{graphicx}
\usepackage{circuitikz}
\usepackage{url}
\usepackage{hyperref}


\begin{document}
\title{PH3244\\Experiment - 6\\Ultrasonic Diffraction}
\author{Aumshree P. Shah\\20231059}
\altaffiliation{\color{red}aumshree.pinkalbenshah@students.iiserpune.ac.in}
\date{\today}
\begin{abstract}
In this experiment we determine the wavelength and velocity of ultrasonic wave in a liquid. We also find bulk modulus and comporessibility of the given liquid. 
\centering
\end{abstract}
\maketitle
\tableofcontents
\vspace*{\fill}
\pagebreak
\pagebreak
\section{Theory and Procedure}
\subsection{Theory}
	From FIG - 1 we see that the ultrasonic waves generated by the transducer travel down the medium (liquid), gets reflected at the bottom (flat glass
plate) of the cell. The incident and reflected waves interfere and a
stationary / standing wave pattern is formed.

The velocity of ultrasonic waves in a liquid is calculated using the formula:

\begin{equation}
V = \mu \Lambda
\label{EQ1}
\end{equation}

where $\mu$ is the frequency of the crystal oscillator and $\Lambda$ is the
wavelength of waves.

We have

\begin{equation}
\Lambda = \frac{n\lambda}{\sin \theta}
\label{EQ2}
\end{equation}

where $n$ is the order of diffraction, $\lambda$ is the wavelength of the
laser used and $\theta$ is the angle of diffraction.

We can find the angle of diffraction by the equation

\begin{equation}
\theta = \tan^{-1}\left(\frac{D}{L}\right)
\label{EQ3}
\end{equation}

where $D$ is the order length and $L$ is the distance measured from the
crystal oscillator to the detector.

The bulk modulus of the liquid is

\begin{equation}
\beta = \rho V^2
\label{EQ4}
\end{equation}

where $\rho$ is the density of the liquid and $V$ is the velocity of the
ultrasonic wave.

The adiabatic compressibility of a liquid can also be calculated using the relation

\begin{equation}
K = \frac{1}{\rho V^2}
\label{EQ5}
\end{equation}
	\begin{figure}[H]
		\centering
		\includegraphics[width=0.5\linewidth]{fig1}
		\caption{Optical system for observation of diffraction by ultrasonic waves.}
	\end{figure}




\subsection{Procedure}Refer to the manual [\cite{ph3244repo}] for a look at procedure, although the required pages will be attached at the back for sumbition along with the signed report.

\section{Observation}
\subsubsection{For 1st crystal\footnote{The least count is to be interpreted from the significant figures written}}
Here we use the plotting graphs method. We put a graph paper very far away and them mark the points in it to measure the distance.
\begin{align*}
	\lambda &= 650\si{\nano\meter}\\
    L &= 2.86\si{\m}\\
    \mu &= 3.25 \text{ MHz}\\
    \rho &= 998 \text{ Kg/m}^3\\
\end{align*}
\begin{table}[H]
\centering
\begin{tabular}{|c|c|c|}
\hline
n (Index) & Distance (cm) & Type of Distance\\
\hline
1 & 0.9 & peak to peak index fringe\\
2 & 1.7 & peak to peak index fringe\\
3 & 1.3 & From center\\
4 & 1.8 & From center\\
\hline
\end{tabular}
\caption{Fringe readings for the above variables.}
\end{table}

\subsection{For 2nd crystal}
Here we use the screw gauge method in which we composate the small $L$ by taking a more presise measurement using a intensity detector.
\begin{align*}
	\lambda &= 650\si{\nano\meter}\\
    L &= 2.34\si{\m}\\
    \mu &= 3.09 \text{ MHz}\\
    \rho &= 998 \text{ Kg/m}^3\\
\end{align*}
\begin{table}[H]
\centering
\begin{tabular}{|c|c|}
\hline
n (Index) & Distance (cm) \\
\hline
1 & 0.9 \\
2 & 1.7 \\
\hline
\end{tabular}
\caption{Fringe readings for the above variables using differnet method.}
\end{table}




\begin{center}
	Data taken on 11 Mar 2026.
\end{center}

\section{Calculations, Analysis and Characteristics}
\subsection{Sources of Error}

\noindent % Prevents paragraph indentation before the first minipage
\begin{minipage}[t]{0.48\textwidth}
\subsubsection{Systematic Errors}
\begin{enumerate}
	\item In alignment of the beam with the crystal.
	\item Improper standing wave produced than the theoratical.
	\item Non Parallel Screen causing increase in distance.
	\item Index of fringe being differnet than what is measured. 
\end{enumerate}
\end{minipage}
\hfill % Distributes remaining horizontal space
\begin{minipage}[t]{0.48\textwidth}
\subsubsection{Random Errors}
\begin{enumerate}
	\item Human errros in measurement.
	\item Uncertanity in instruments.
	\item Shaking of table resulting in moving fringes.
	\item Diffraction increasing the fringe width.
\end{enumerate}
\end{minipage}

\subsection{Error Propogation}
To calculate the standard deviation from values we have measured with some least count, and the values we have applied digitally, we use the formula [\cite{tewari_uncertainty}]: 

\begin{equation}
\sigma = \frac{\text{Least Count}}{\sqrt{12}}
\end{equation}
and to calculate the standard deviation for values we have applied on analog instrument with some least count, we use the formula [\cite{tewari_uncertainty}]: 
\begin{equation}
\sigma = \frac{\text{Least Count}}{\sqrt{24}}
\end{equation}
Since all our observations only satisfy the criteria for Equation-6 we will be using that. Now from Equation-1,2,3 we can write $V$ as : 
\begin{equation}
	V = \mu n\lambda \frac{\sqrt{{D^2+L^2}}}{D}
\end{equation}
to calculate the error of $V$ we notice that it'll only (my only here I mean the most significant factor of error) will be the uncertainty in $D$, hence in error calculation we will assume other errors are negligible.
Now using formula for error propogaiton [\cite{taylor_error_analysis}], efor Equation-8 we see that: 
\begin{align*}
	\sigma_V / \sigma_D &=  \frac {\text d}{\text d D} V\\
		 &= \mu n \lambda \frac 1 D\frac {\text d}{\text d D}  \left( \sqrt{D^2+L^2} \right) + \mu n \lambda \sqrt{D^2+L^2}\frac {\text d}{\text d D}  \left( \frac 1 D \right) 
\end{align*}
\begin{equation}
	\implies \sigma_V	 =\boxed{\left( \mu n \lambda \frac {1}{\sqrt{D^2+L^2}} - \mu n \lambda \sqrt{D^2+L^2}\frac {1}{D^2} \right) \sigma_D}
\end{equation}
Similarly from Eqation-4,5 we get the uncertainty in $\beta$ and $K$ as:
\begin{equation}
	\sigma_\beta = 2V\rho \sigma_V
\end{equation}
\begin{equation}
	\sigma_K = \frac {2}{\rho V^3} \sigma_V
\end{equation}

\subsection{Calculations}

First we see that since the Table-2 has only 2 data points, where also there is a good chance that it might have high precision and low accuract, we wil not consider the data from Table-2 in our calculations. Using Equation-8,4,5 to calculate the variables and Equation-9,10,11 to calcualte their uncertainty, we get:
\begin{table}[H]
\centering
\begin{tabular}{c c c c c c c}
\hline
$n$ & Distance & Least Count & SI Distance & SI Uncertainty & Velocity & Velocity Uncertainty  \\
    & (cm) & (cm) & (m) & (m) & (m/sec) & (m/sec) \\
\hline
1 & 0.45 & 0.05 & 0.0045 & 0.00014 & 1343 & 42 \\
2 & 0.85 & 0.05 & 0.0085 & 0.00014 & 1422 & 23 \\
3 & 1.30 & 0.10 & 0.0130 & 0.00029 & 1394 & 29 \\
4 & 1.80 & 0.10 & 0.0180 & 0.00029 & 1343 & 19 \\
\hline
\end{tabular}

\caption{Calculating velocity from TABLE-1}
\end{table}
Using the weighted average formula [\cite{standard_error_wiki}] we can calculate $V$, $\beta$ and $K$ which will be discussed in results. 
\begin{equation}
	\bar{v}=\frac{\sum_{i=1}^{N} \frac{v_i}{\sigma_i^2}}{\sum_{i=1}^{N} \frac{1}{\sigma_i^2}}
\end{equation}
\begin{equation}
	\sigma_{\bar{v}} = \sqrt{\frac{1}{\sum_{i=1}^{N} \frac{1}{\sigma_i^2}}}
\end{equation}
e weighted average we get is: From this 

\section{Result and Conclusion}
By taking the weighted average of velocity and Equaiton 12,13 of each value from TABLE-3, we get value of $V$ as $\boxed{V = 1375 \pm 12 \text{ m/sec}}$; From Equation-4,10 we get the bulk modulus as $\boxed{\beta = 1.88 \pm 0.03 \text{ GPa}}$ and finally from Equation-5,11 we get the compressiblity as $\boxed{K = 53.67 \pm 0.08 \times 10^{-6}\text{ atm}}$ . We will compare this values with the theoratical values in Appendix-A.

\vspace{1cm}
\hrule
\vspace{1cm}
\appendix

\bibliographystyle{apsrev4-2}
\bibliography{refrences}
\hfill\pagebreak

\section{Theoratical values and comparision}
The theoratical value for bulk modulus of water at 20\si{\celsius} is around 2.2 GPa, this is 16\% different from the calculated while the uncertainty being only 1\%. The reason for this are the systematic errors discussed in the error analysis section.
\section{Improvments in procedure}
We can measure the $D$ on a flourosent paper in dark room which will make the readings a lot easier and accurte. Furthermore we can calculate the $D$ at very large $L$ this will make the relative uncertainty of $D$ small and thus giving more accurate readings. For the random error, we can take more readings to reduce them.
\end{document}
